\chapter{KẾT LUẬN VÀ HƯỚNG PHÁT TRIỂN}
\label{chap:ketluan}

\section{Kết luận}
\label{sec:ketluan}

Dự án đã hoàn thành đầy đủ các mục tiêu nghiên cứu đề ra ở mục~\ref{subsec:muctieu}, với các kết quả cụ thể sau:
\begin{enumerate}
    \item Nghiên cứu và hệ thống hóa toàn diện cơ sở lý thuyết về Lọc cộng tác, Nhân tử hóa ma trận (MF), Generalized Matrix Factorization (GMF), Multi-Layer Perceptron (MLP), kiến trúc mô hình lai Neural Matrix Factorization (NeuMF) \parencite{he2017ncf, koren2009matrix}, đồng thời khảo sát các hướng nghiên cứu hiện đại hơn (Graph-based CF, Sequential Recommendation, LLM/Diffusion-based Recommendation, mục~\ref{sec:modern_trends}) để xác định rõ vị trí và giới hạn của hướng tiếp cận đã chọn.
    \item Xây dựng hoàn chỉnh quy trình tiền xử lý dữ liệu, áp dụng \textbf{đồng nhất trên hai bộ dữ liệu độc lập} có đặc trưng cấu trúc khác biệt rõ rệt: DataCo Smart Supply Chain \parencite{dataco2019} (catalog nhỏ, dày đặc) và H\&M Personalized Fashion Recommendations (catalog lớn, cực thưa) -- bao gồm lọc $k$-core, ánh xạ ID, chuyển đổi phản hồi ẩn nhị phân, phân chia Leave-One-Out (kèm kiểm tra Disjoint bắt buộc), và một giải pháp kỹ thuật tiền xử lý cache Parquet giúp xử lý tệp giao dịch H\&M 31,8 triệu dòng (3,49GB) trên phần cứng CPU phổ thông (mục~\ref{subsec:hm_caching}).
    \item Thiết kế, cài đặt và huấn luyện thành công các mô hình GMF, MLP, Early Fusion và NeuMF (from-scratch \& pre-trained) bằng framework PyTorch \parencite{paszke2019pytorch}, kết hợp kỹ thuật lấy mẫu tiêu cực động và chiến lược tiền huấn luyện.
    \item Kiểm chứng thực nghiệm định lượng, trung thực trên cả hai bộ dữ liệu (Chương~\ref{chap:thucnghiem}): NeuMF vượt trội rõ rệt các baseline cổ điển không dùng mạng nơ-ron (Random, Most Popular, ItemKNN, BPR-MF) trên cả hai giao thức Full Ranking và Sampled-99, nhưng \textbf{chưa cho thấy cải thiện rõ rệt} so với chính các nhánh thành phần của nó (GMF, MLP) khi đứng riêng -- một phát hiện khách quan, không như kỳ vọng ban đầu của khung lý thuyết NCF, được phân tích chi tiết ở mục~\ref{subsec:ablation}. Đồng thời, thực nghiệm cho thấy lựa chọn Early Fusion hay Late Fusion (NeuMF) tối ưu phụ thuộc vào đặc trưng mật độ/quy mô catalog của dữ liệu, không có câu trả lời cố định.
    \item Hoàn thành phân tích bóc tách thành phần (Ablation Study) và đánh giá phân tầng Popular/Long-tail, phát hiện dấu hiệu thiên vị phổ biến (\textit{popularity bias}) rõ rệt ở nhiều mô hình, đặc biệt trên H\&M.
\end{enumerate}

\section{Hạn chế của đề tài}
\label{sec:hanche}

\begin{itemize}
    \item Mô hình NeuMF hiện tại chủ yếu khai thác ma trận tương tác User--Item dựa trên mã ID và giá trị đơn hàng, chưa tích hợp các đặc trưng đa phương thức (Multimodal Information) như hình ảnh hoặc mô tả văn bản chi tiết của sản phẩm.
    \item Chưa mô hình hóa tính động và sự thay đổi sở thích ngắn hạn của người dùng theo chuỗi thời gian (Sequential/Session-based Recommendation).
    \item Bài toán khởi đầu lạnh (Cold-start) cho người dùng/sản phẩm hoàn toàn mới \textbf{nằm ngoài phạm vi thực nghiệm} của dự án -- mô hình hiện tại chỉ hoạt động với các thực thể đã có đủ lịch sử tương tác sau khi lọc $k$-core.
    \item \textbf{Chưa kiểm định ý nghĩa thống kê đa seed:} Kết quả ở Chương~\ref{chap:thucnghiem} được báo cáo trên một seed duy nhất (seed=42) cho mỗi bộ dữ liệu (mục~\ref{subsec:multiseed_results}); do đó, chênh lệch nhỏ giữa NeuMF và các nhánh thành phần \textbf{chưa được xác nhận là có ý nghĩa thống kê} hay chỉ là nhiễu ngẫu nhiên của một lần chạy -- đây là hạn chế quan trọng nhất cần khắc phục trước khi công bố kết luận chính thức.
    \item \textbf{Kết quả H\&M mới ở quy mô lát cắt:} Thực nghiệm định lượng trên H\&M được thực hiện trên lát cắt 100.000 dòng giao dịch đầu (1.743 người dùng, 1.080 sản phẩm), không phải toàn bộ 31,8 triệu dòng gốc (889.062 người dùng, 90.690 sản phẩm ở quy mô đầy đủ, mục~\ref{subsec:hm_apply}) do chi phí tính toán của giao thức Full Ranking tăng theo tích số người dùng $\times$ số sản phẩm; các kết luận rút ra từ H\&M cần được xác nhận lại ở quy mô đầy đủ.
\end{itemize}

\section{Hướng phát triển}
\label{sec:huongphattrien}

\begin{enumerate}
    \item Mở rộng sang kiến trúc Mạng nơ-ron đồ thị (Graph Neural Networks -- GNNs) như LightGCN nhằm khai thác tương tác bậc cao giữa User và Item trong không gian thưa dữ liệu.
    \item Ứng dụng cơ chế Self-Attention (SASRec/Transformer) để mô hình hóa hành vi mua sắm tuần tự theo thời gian thực của khách hàng.
    \item Tích hợp mô hình ngôn ngữ lớn (LLMs) và mô hình thị giác máy tính để trích xuất đặc trưng đa phương thức của sản phẩm, giải quyết triệt để bài toán Cold-start.
    \item Mở rộng mô hình theo hướng học đa tác vụ (Multi-Task Learning) nhằm dự đoán đồng thời các hành vi nhấp chuột (Click), thêm vào giỏ hàng (Cart) và mua sắm (Purchase).
    \item Mở rộng thực nghiệm H\&M lên toàn bộ quy mô dữ liệu (889.062 người dùng, 90.690 sản phẩm sau lọc $k$-core=5, mục~\ref{subsec:hm_apply}) bằng cách chuyển giao thức đánh giá chính sang Sampled-99 (chi phí không phụ thuộc kích thước catalog) thay vì Full Ranking, đồng thời chạy Full Ranking như một phép đối chiếu trên mẫu con người dùng được lấy ngẫu nhiên định kỳ; đây là điều kiện tiên quyết để kiểm chứng các kết luận ở Chương~\ref{chap:thucnghiem} (rút ra từ lát cắt 100.000 dòng) có còn đúng ở quy mô đầy đủ hay không.
    \item Hoàn tất kiểm định ý nghĩa thống kê đa seed (mục~\ref{subsec:multiseed_results}) trên cả hai bộ dữ liệu -- hạ tầng đã sẵn sàng (\texttt{scripts/04\_multi\_seed.py}, \texttt{scripts/06\_aggregate\_seeds.py}) nhưng chưa được chạy đầy đủ do giới hạn thời gian của lần kiểm tra số liệu này.
\end{enumerate}

\section{Triển khai thực tế giao diện thực nghiệm}
\label{sec:huongthietke_giaodien}

Thiết kế đầy đủ của hệ thống giao diện thực nghiệm (kiến trúc ba tầng, đặc tả API, luồng xử lý, lựa chọn mô hình triển khai) đã được trình bày như một cấu phần bắt buộc của đề tài ở mục~\ref{sec:demo_design}, không phải một hướng mở rộng tùy chọn. Phần còn lại -- \textbf{cài đặt và triển khai thực tế} bản thiết kế đó thành mã nguồn chạy được -- là công việc tiếp theo ngay sau báo cáo này, với các thành phần đầu vào đã sẵn sàng:
\begin{itemize}
    \item \textbf{Mô hình đã huấn luyện:} Trọng số của mọi mô hình (GMF, MLP, Early Fusion, NeuMF from-scratch/pre-trained) trên cả hai bộ dữ liệu đã được lưu tại \texttt{outputs/checkpoints/} dưới định dạng PyTorch (\texttt{.pt}), đúng như Tầng 3 của thiết kế (mục~\ref{subsec:demo_architecture}) yêu cầu.
    \item \textbf{Bảng ánh xạ định danh:} hai tệp \texttt{user\_id\_map} và \texttt{item\_id\_map} (định dạng Parquet, mục~\ref{subsec:hm_caching}), phục vụ trực tiếp bước 1 và bước 5 trong luồng xử lý của endpoint \texttt{/recommend} (mục~\ref{subsec:demo_api}).
    \item \textbf{Việc cần làm tiếp theo:} cài đặt mã nguồn thực tế cho Tầng 1 (giao diện) và Tầng 2 (API suy diễn) theo đúng đặc tả ở Bảng~\ref{tab:demo_api}, đóng gói bằng Docker, và đo đạc thực tế thời gian phản hồi để xác nhận mục tiêu phi chức năng đã đặt ra ở mục~\ref{subsec:demo_model_choice}.
\end{itemize}